\section{Cross-experiment analysis}\label{sec:analysis}

The experiments of \cref{sec:big,sec:small,sec:case,sec:programmes} share one task pool, one runner and, for E4B, one
baseline. This section reads across them. It asks how large the noise of each kind of comparison is, which tasks can
move at all, what a hosted score adds to a local result, and what the record implies for the test of a further
candidate. Every analysis in this section was chosen after the results were known. None of them changes a registered
verdict.

\subsection{Three comparisons and their noise}\label{sec:noise}

\begin{table}[t]
\centering
\caption{Standard deviation of the difference of two arms that run on the same tasks, from run-to-run variation
alone. The first two rows are estimated from the repeated baseline sessions (\cref{sec:small}); the third row uses
the repeated public draws of unchanged 31B files \citep{afshar2026census}.}
\label{tab:noise}
\footnotesize
\begin{tabular}{lrrl}
\toprule
Comparison of two arms & Runs per arm & SD, points & Time as run \\
\midrule
Screen on S (48 tasks, 2 repeats) & 96 & \ScreenSdDiffPts{} & 1--2 hours (3--6 arms) \\
Confirmation on H (128 tasks, 4 repeats) & 512 & \NoiseSdDiffPts{} & about 3 hours (2 arms) \\
Hosted, one draw per file (58 public tasks) & 58 & \HostedSdDiffPts{} & 2 days \\
\bottomrule
\end{tabular}
\end{table}

\Cref{tab:noise} puts the three comparisons that we used side by side. All three hold the tasks fixed, so the numbers
answer one question only: would the difference appear again if the same two builds ran again on the same tasks? For
that question a confirmation block on the workstation is the most exact instrument that we have, and a pair of hosted
scores is the least exact.

The pass rule asks a harder question: would the gain hold on tasks from repositories that the block did not contain?
Its test resamples repositories, and its answer can differ widely from the answer for fixed tasks. The third
application shows the size of that difference. The candidate solved \CandThreeDiffRuns{} more runs than the baseline,
a gain of $\GainThree$ points. Against run-to-run noise alone this is \CandThreeNoiseRatio{} standard deviations.
Under the repository-clustered test the one-sided $p$ is \GatePThree{}, above the limit of 0.025
(\cref{sec:leadgain}). Both statements are correct. The first says that the two builds very probably differ on the
128 tasks of H. The second says that the 19 repositories of H do not yet show that the gain carries over to other
repositories.

A third source of error is selection. The candidate that enters a confirmation is the leader of a screen, and the
leader of several equal arms is ahead by construction: \ScreenNullLeadFive{} points on average when five variants are
screened against the baseline, and \ScreenNullLeadTwo{} points when two are (\cref{sec:leadgain}). A reader who compares a screen
lead with \cref{tab:noise} must subtract this first.

\subsection{Tasks that the small model does not solve}\label{sec:never}

On most tasks the outcome of an E4B run does not depend on the run, and it does not depend on the arm. The baseline
solved \HNever{} of the 128 tasks of H in none of its eight runs. Each of these tasks was also run eight times by the
two candidates that were tested on H. They solved \HNeverRescued{} of them at least once, in \HNeverRescuedRuns{} of \HNeverCandRuns{}
runs; \HNeverAny{} tasks stayed unsolved in all 16 runs. On the screen set S the baseline never solved \SNeverBase{}
of 48 tasks in six runs, and \SNeverAny{} tasks stayed unsolved in the 24 runs of all arms of the three screens. Over
all E4B runs on the workstation, \WsTasksAny{} tasks were attempted, each in \WsRunsPerTaskMin{} to
\WsRunsPerTaskMax{} runs. \WsTasksSolvedAny{} of them were solved at least once. The other
\WsTasksNeverPct\% were never solved.

``Never'' is a statement about our runs. A task that is solved in one run of ten stays unsolved in eight runs with a
probability of 43\%. The few candidate runs that solved a task of the never-solved group are therefore what such tasks
would give without any change to the agent, and we do not read them as tasks that a candidate opened.

Two consequences follow. First, the changes that we tested act on a minority of tasks. Even the largest gain,
that of the third candidate, changed the solved count of \CandThreeTasksUp{} tasks upwards and of \CandThreeTasksDown{}
tasks downwards, and left \CandThreeTasksSame{} of 128 as they were. Second, the task set fixes most of a solved share before any arm
runs. The share of tasks that some arm solved in some run, \WsTasksSolvedAny{} of \WsTasksAny{}, is about twice the
baseline's solved share of runs (\BaseWsPct\%). It bounds what any choice among our variants could reach on these
tasks. \Cref{sec:cost} showed how the failures look: in \WsUnsolvedWrongPct\% of the unsolved workstation runs the
agent handed in a patch that the tests reject, and in \WsUnsolvedTimeoutPct\% it was stopped by the time limit.

\subsection{What a hosted score adds}\label{sec:hostednoise}

\paragraph{Precision.}
The census of repeated public draws gives a standard deviation of \HostedSdTasks{} tasks of 58 for one draw of
unchanged 31B files (our recount of its 29 draws of six file sets). The difference of two submissions then has a
standard deviation of \HostedSdDiffPts{} points. To detect a gain of two tasks of 58 with a power of 80\% at a
one-sided level of 0.025, each of the two files would need about \HostedDrawsNeeded{} draws. At one submission per
day this is twenty days for one comparison. The difference of two arms in one confirmation block has about
\LocalBlockRatio{} times less run-to-run variance than the difference of two single hosted draws, and the block takes
three hours. For E4B files no repeated draws exist, so the first number is not known for the small model.

\paragraph{Relevance.}
The hosted score is measured on the tasks that decide the competition. Local runs are not. For the four arms with at
least 70 local runs, the hosted score is between \HostedRatioMin{} and \HostedRatioMax{} of the local solved share
(\cref{tab:hosted}): 0.12 and 0.18 against 54\% and 53\% of the development tasks for the two 31B arms, and 0.06
and 0.05 against 20\% and 23\% of Berth runs for the two E4B arms. The hidden tasks need not resemble either local
set. The development tasks are public, and the Berth tasks are those that load and pass validation. We did not test
which of these reasons accounts for the gap.

\paragraph{Order.}
We have two pairs of arms with both a local and a hosted result. \arm{W3} and \arm{BF5} tied locally (38 and 37 of
70) and scored 0.12 and 0.18. The third E4B candidate was ahead of the baseline locally by $\GainThree$ points and
scored 0.05 against 0.06. In both pairs the hosted order is the reverse of the local order. The hosted difference of
the 31B pair is about one and a half standard deviations of a difference of single draws (\cref{sec:hosted}); that
of the E4B pair is 0.01. Neither can be told from noise. Two pairs cannot show whether a local gain carries over to
the hidden tasks. They show that one hosted draw per arm cannot answer the question.

\subsection{Time limits and the two models}\label{sec:limits-models}

The time limit of a session shapes the results of the large model and hardly touches the small one. On the development
tasks under emulated latency, \TimeoutBigDevPct\% of the 31B runs were stopped by the limit; on Berth tasks on
4$\times$L4, \TimeoutBigBerthPct\%. Of the E4B runs on the workstation, \TimeoutSmallWsPct\% were stopped by the time
limit and \BudgetSmallWsPct\% by the limit on tool calls (\cref{tab:cost}). For 31B the limit binds in about one
run of eight, so the way an arm uses its time is part of its result. The first wrapped arm, for example, finished its
jobs sooner than the reference because it wrote no history summaries (\cref{sec:big-arms}). For E4B, time is not the
binding limit. The three candidates used \RecheckTimeOne{}, \RecheckTimeTwo{} and \RecheckTimeThree{} times the
agent time of the baseline, inside the factor of 1.2 that the pass rule allows. With so few runs stopped by a limit,
a faster E4B agent could gain at most that fraction of runs. A change that helps must alter the patch that the agent
submits.

\subsection{What the record implies for a further candidate}\label{sec:next}

\begin{enumerate}[leftmargin=*,itemsep=2pt,topsep=3pt]
\item \textbf{Confirmations are the scarce resource.} A screen costs an hour and ranks nothing
(\cref{sec:leadgain}). A confirmation costs about a thousand runs and three hours, and it spends its panel: H has
been used twice, and the third candidate was built after the first result on H was known.
\item \textbf{The validated pool is nearly used up.} Of the \PoolValid{} validated tasks, \PoolUsed{} have been
used by some experiment of this paper. The remaining \PoolUntouched{} tasks lie in \PoolUntouchedRepos{}
repositories; one repository holds \PoolUntouchedTopRepo{} of them and \PoolUntouchedSmallRepos{} repositories hold
one task each. The pass rule asks for three repositories with a gain, so these tasks alone cannot carry a
confirmation. A further \CaseOnlyTasks{} tasks were used only by the case study of \cref{sec:case}, which compared
prompt renderings and no candidate. They can serve as a fresh panel only for a candidate whose design did not draw
on the per-task results of that study.
\item \textbf{The size of the panel sets the power.} For an even gain of three tasks of 58 the simulated pass rate
of the rule is about 60\% on the small panel and about 90\% on the large one (\cref{tab:simulated}). A gain that is
concentrated in few repositories passes far less often.
\item \textbf{A hosted draw is a check of the file.} It shows that a bundle loads and runs under the scorer. It does
not show which of two bundles is better (\cref{sec:hostednoise}).
\end{enumerate}
